\documentclass[11pt]{article}
\usepackage{acl}
\usepackage{times}
\usepackage{latexsym}
\usepackage{booktabs}
\usepackage{amsmath}
\usepackage{microtype}
\usepackage{graphicx}
\usepackage{url}

\title{A Validity-First LoRA Baseline for Text-to-SVG Generation}

\author{
Siva Srinivas Venigalla \\
New York University Tandon School of Engineering \\
\texttt{sv2881@nyu.edu}
}

\begin{document}
\maketitle

\begin{abstract}
This report presents a reproducible baseline for the DL Spring 2026 text-to-SVG Kaggle competition. The system fine-tunes \texttt{Qwen2.5-1.5B-Instruct} with LoRA on the provided prompt--SVG pairs, then applies strict SVG repair, validation, and a TF-IDF retrieval fallback to satisfy competition constraints. The saved notebook run demonstrates a 100\% syntactic validity rate on the generated submission file and completes inference for 1,000 hidden prompts in 57.32 minutes on a Tesla T4 GPU. However, qualitative inspection shows that the current baseline often collapses to a generic fallback SVG, indicating that validity is high but semantic fidelity remains weak. This makes the current system a useful reproducible baseline, but not yet a competitive final entry.
\end{abstract}

\section{Introduction}
The DL Spring 2026 Kaggle contest asks students to generate valid SVG code directly from natural-language prompts. Unlike many open-ended generative tasks, this competition applies a hard validity gate before scoring semantic quality. Any SVG that fails parsing, violates the allowed-tag set, exceeds complexity limits, or fails to render receives a zero for that sample. As a result, a reasonable first milestone is to build a pipeline that is reliably valid under Kaggle's execution and safety constraints.

The baseline in this repository is designed around that objective. It uses a compact instruction-tuned language model, parameter-efficient adaptation with LoRA \citep{hu2022lora}, quantized loading in the style of QLoRA \citep{dettmers2023qlora}, and a lightweight retrieval fallback to recover from malformed generations. The central design choice is pragmatic: prioritize a stable, reproducible, Kaggle-compatible system before optimizing leaderboard performance.

\section{Dataset}
The notebook uses only the provided competition data. In the saved run, the training file contains 50,000 prompt--SVG pairs and the hidden-test file contains 1,000 prompts. Each training row provides an \texttt{id}, a natural-language \texttt{prompt}, and a target \texttt{svg}. The test set contains \texttt{id} and \texttt{prompt} only.

Before training, the notebook performs lightweight filtering and repair:

\begin{itemize}
\item empty prompts are removed;
\item SVG strings are normalized by extracting the root \texttt{<svg>} block when present;
\item missing namespace or canvas attributes are inserted;
\item XML declarations, doctypes, and comments are removed;
\item overlong or malformed SVGs are rejected.
\end{itemize}

All 50,000 rows in the saved run passed the initial cleaning stage. For feasibility on modest GPU hardware, the notebook then constructs a balanced subset by binning prompt lengths and SVG lengths, producing 8,500 usable examples. From this subset, 8,000 rows are used for training and 425 for validation.

\section{Methodology}
\subsection{Model}
The base model is \texttt{Qwen/Qwen2.5-1.5B-Instruct} \citep{qwen25}. This choice keeps the parameter count below the course guidance while preserving enough capacity to model the mapping from text descriptions to short SVG programs.

LoRA adapters are attached to the projection layers
\texttt{q\_proj}, \texttt{k\_proj}, \texttt{v\_proj}, \texttt{o\_proj},
\texttt{up\_proj}, \texttt{down\_proj}, and \texttt{gate\_proj}.
The saved run reports 18,464,768 trainable parameters out of 1,562,179,072 total parameters, or roughly 1.18\% of the model.

\subsection{Prompting Format}
Each example is converted into a short instruction-style text block:
\begin{quote}
\small
\texttt{System: Generate a compact valid SVG ...}\\
\texttt{User: <prompt>}\\
\texttt{Assistant: <svg>}
\end{quote}

This formatting keeps the task explicit while remaining simple enough for fast tokenization and fine-tuning. During inference, the same prompt structure is used but the assistant response is left blank for generation.

\subsection{Quantization and Training Strategy}
When a GPU is available, the notebook attempts 4-bit NF4 quantized loading with double quantization and half-precision compute. This substantially reduces memory usage and allows LoRA fine-tuning on a Tesla T4 GPU. If 4-bit loading fails, the notebook falls back to a standard model load.

\subsection{SVG Repair and Validation}
Because invalid SVGs receive zero credit, the pipeline includes several post-processing checks:

\begin{itemize}
\item extraction of the generated SVG span from raw text;
\item reinsertion of required \texttt{xmlns}, \texttt{width}, \texttt{height}, and \texttt{viewBox} attributes when missing;
\item length filtering against the 16,000-character limit;
\item XML parsing with \texttt{lxml} to reject malformed outputs.
\end{itemize}

If generation fails these checks, the notebook falls back to a retrieval-based SVG or, in the worst case, to a minimal hand-written SVG consisting of a white background and a black circle. This fallback is intentionally simple: it protects the submission from invalid outputs, but may sacrifice semantic quality.

\subsection{Retrieval Fallback}
The notebook includes a TF-IDF prompt retriever built over the cleaned training prompts. For a malformed generation, the system retrieves the nearest prompt in the training subset and reuses the corresponding cleaned SVG. This component is not the main modeling strategy; it is included as a practical safeguard for Kaggle submission robustness.

\section{Experimental Setup}
Table~\ref{tab:hparams} lists the main hyperparameters used in the saved notebook run.

\begin{table}[t]
\centering
\small
\begin{tabular}{ll}
\toprule
Setting & Value \\
\midrule
Base model & \texttt{Qwen2.5-1.5B-Instruct} \\
Training rows & 8,000 \\
Validation rows & 425 \\
Sequence length & 512 \\
Max new tokens & 180 \\
LoRA rank & 16 \\
LoRA alpha & 16 \\
LoRA dropout & 0.05 \\
Epochs & 1 \\
Per-device batch size & 1 \\
Gradient accumulation & 4 \\
Learning rate & $2 \times 10^{-4}$ \\
Warmup steps & 30 \\
Weight decay & 0.01 \\
Inference batch size & 8 \\
Seed & 42 \\
\bottomrule
\end{tabular}
\caption{Key hyperparameters used in the baseline notebook.}
\label{tab:hparams}
\end{table}

The saved notebook output reports a Tesla T4 GPU and successful 4-bit model loading. Inference over the 1,000 hidden test prompts completed in 57.32 minutes, which fits the intended notebook-based workflow but leaves limited budget for larger models or expensive search.

\section{Results}
Table~\ref{tab:results} summarizes the directly observed outcomes from the saved notebook execution.

\begin{table}[t]
\centering
\small
\begin{tabular}{ll}
\toprule
Metric & Observed value \\
\midrule
Submission rows & 1,000 \\
Rows starting with \texttt{<svg>} & 100.0\% \\
Valid SVG rate & 100.0\% \\
Inference time & 57.32 minutes \\
Saved Kaggle format & Yes \\
Public leaderboard score & \textbf{TBD} \\
Private leaderboard score & \textbf{TBD} \\
\bottomrule
\end{tabular}
\caption{Results directly supported by the notebook outputs.}
\label{tab:results}
\end{table}

The strongest outcome of this baseline is validity: every row in the final submission passes the notebook's SVG validity check, and every row begins with a proper \texttt{<svg>} root tag. This directly supports the design goal of avoiding zero-score samples caused by malformed outputs.

The main weakness is semantic collapse. The notebook's sanity-check examples repeatedly produce the same fallback SVG instead of prompt-specific graphics. This suggests that the current model either fails to generate useful SVGs after fine-tuning or is triggering the fallback path too frequently. In either case, the observed behavior indicates low visual fidelity despite perfect validity.

\section{Ablations and Failure Analysis}
The current repository contains one end-to-end baseline configuration rather than a full ablation sweep. That is a limitation of the present report. For the final version, the following controlled ablations should be added because they directly test the assumptions of the pipeline:

\begin{itemize}
\item \textbf{No retrieval fallback}: measure how much of the 100\% validity rate comes from the backup mechanism rather than the generative model itself.
\item \textbf{No SVG repair}: measure how often the raw model output is already valid before post-processing.
\item \textbf{Different training subset sizes}: compare 2k, 8k, and full-data fine-tuning to test whether the model is underfit.
\item \textbf{Different LoRA ranks}: compare rank 8, 16, and 32 to test capacity limits under the same compute budget.
\item \textbf{Different decoding limits}: compare max generation lengths and decoding strategies to see whether truncation is contributing to malformed or generic outputs.
\end{itemize}

Even without the full sweep, the saved outputs already reveal one clear failure mode: the validity-oriented fallback dominates qualitative behavior. This means the current baseline is conservative and robust, but it is not yet learning prompt-conditioned SVG structure well enough to compete on the leaderboard.

\section{Code and Reproducibility}
The code repository for this project is:
\begin{quote}
\small
\url{TODO_REPO_URL}
\end{quote}

The notebook fixes random seeds for Python, NumPy, and PyTorch and produces a Kaggle-ready CSV file with the required \texttt{id,svg} schema. The repository also includes a dependency list in \texttt{requirements.txt} and preserves the exact notebook used to generate the baseline results.

The report should also include a public link to the final trained adapter weights. At the time of writing, that link still needs to be added:
\begin{quote}
\small
\texttt{[Add model-weights URL here]}
\end{quote}

\paragraph{AI Tool Disclosure}
LLM-based tooling was used for code cleanup, debugging assistance, repository organization, and report drafting. Final methodological decisions, experiment selection, and result interpretation were manually reviewed by the author.

\section{Conclusion}
This project establishes a reproducible baseline for the DL Spring 2026 text-to-SVG competition. The current system is strong on validity and Kaggle compatibility: it fine-tunes within modest hardware limits, generates a fully valid submission file, and avoids malformed SVG outputs. However, the saved qualitative outputs show that the model is not yet producing prompt-specific graphics reliably, so leaderboard performance is likely limited in its current form.

The next stage should focus on reducing dependence on the fallback path and improving prompt-conditioned SVG generation quality. The most important near-term work is a controlled ablation study over retrieval fallback, repair logic, training subset size, and LoRA capacity, followed by measurement on the Kaggle public leaderboard.

\bibliographystyle{acl_natbib}
\bibliography{references}

\end{document}
